\documentclass{article}

\usepackage{booktabs}
\usepackage{tabularx}
\usepackage{tabularx}
\usepackage{array}
\usepackage{float}
\usepackage{hyperref}

\title{Development Plan\\\progname}

\author{\authname}

\date{}

\input{../Comments.text}
\input{../Common.text}

\begin{document}

\maketitle

\begin{table}[hp]
\caption{Revision History} \label{TblRevisionHistory}
\begin{tabularx}{\textwidth}{llX}
\toprule
\textbf{Date} & \textbf{Developer(s)} & \textbf{Change}\\
\midrule
September 27 & Nathan & Changed project decomposition and scheduling, changed GitHub Projects usage, changed coding standards, changed Appendix -- Team Charter\\
September 27 & Aidan & Added sections 5 to 7 of the development plan regarding Team communication plan, roles, and workflow. \\
September 28 & Robert & Added sections 1-4 and 9-10\\
\bottomrule
\end{tabularx}
\end{table}


\newpage{}

This document outlines the processes, tools, and practices that the team will use to plan, develop, and manage the \href{https://github.com/robertLam04/daiim}{DAIIM} project. It describes the team's approach to communication, collaboration, project scheduling, workflow, and software development, as well as the technologies and coding standards that will guide implementation. The plan also identifies the project's proof-of-concept strategy and describes how generative AI and other development tools will be incorporated into the development process.

\section{Confidential Information?}
This project does not contain any confidential or protected information.
\section{IP to Protect}

The team does not have any project-specific intellectual property that requires protection through an IP agreement. The DAIIM source code will be released as open-source software under the license described in the copyright license section below.

\section{Copyright License}

This project, including all source code and documentation, will be under the MIT license. The \href{https://github.com/robertLam04/daiim/blob/main/LICENSE}{license} can be found in the project repository.

\section{Team Meeting Plan}

\begin{table}[hp]
\caption{Team and Supervisor Meeting Schedule} \label{TblMeetingSchedule}
\begin{tabularx}{\textwidth}{lXXX}
\toprule
\textbf{Meeting} & \textbf{Location} & \textbf{Frequency} & \textbf{Time}\\
\midrule
Team meeting & McMaster campus & Weekly & Wednesday 1:30 pm\\
Supervisor meeting & Microsoft Teams & Biweekly & TBD\\
\bottomrule
\end{tabularx}
\end{table}

All meetings will be recorded in a Github issue outlining the date, attendance and agenda. Each meeting will have a chair who is responsible for creating the agenda and sharing it with others prior to the meeting. If possible the agenda should be shared 12hrs in advance to give team members a chance to review and propose different topics. For each meeting, the chair will ask for a volunteer from the team to record meeting minutes. If no one volunteers, the note-taking role will rotate among team members. Meeting minutes and decisions made during the meeting will be recorded as a comment on the meeting issue. Ad hoc meetings may also be scheduled as necessary for urgent decisions, issue
resolution, or preparation for a deliverable.

\section{Team Communication Plan}

The team will use a Microsoft Teams group chat for regular communication. 
The team also has a shared Word document for collaborative editing of documents and planning during meetings.  Github issues will be used to track each meeting and any associated agendas/minutes. \newline
\noindent Communication with the supervisor will be done through email and Microsoft Teams. 
Every email with the supervisor will be sent to all team members. 
The team will meet regularly with the supervisor, subject to the supervisor's availability.  \newline

\noindent Github issues will be the primary method of communication for development tasks. Each issue will represent a development task and will be assigned to a team member; the team member will be responsible for updating the issue with their progress.  
The issues will be associated with a milestone representing the deliverable, and will be placed appropriately in the Kanban project according to its status (e.g., To Do, In Progress, Done).
Feedback is expected from at least one other team member for each issue, through either a comment on the Github issue or reviewed in a meeting or group chat. \newline


\section{Team Member Roles}

\begin{itemize}
  \item \textbf{Agenda creator}: The agenda creator is responsible for creating
  the agenda for a meeting, and posting it to the meeting issue. The agenda creator will 
  use the template for team meetings located in the repo, containing the agenda and an attendance tracker. 
  
  \item \textbf{Notetaker}: The notetaker is responsible for taking notes during a meeting, and posting the meeting minutes to the meeting issue. 
  The meeting minutes will be posted as a comment under the meeting issue. 
  Multiple people may contribute to the notetaker role during a meeting, especially if there are multiple topics / individual updates from each person.
  
  \item \textbf{Pull request reviewer}:
  The pull request reviewer role is a role that is assigned to a team member for each pull request. The pull request reviewer is responsible for reviewing the code changes in the pull request, providing feedback, and approving or requesting changes to the pull request.
  Every team member is expected to review pull requests, but the pull request reviewer is the primary reviewer for that pull request. They are in charge of making the final decision on approving the request. 

  \item \textbf{Developer}: The developer role is a role assigned to a team member for each development task. The developer is responsible for implementing the task, and updating the associated issue with their progress.
  Every team member is expected to take on this role and contribute to development tasks as necessary, but the developer is the primary person responsible for completing a given task.
  
\end{itemize}


\section{Workflow Plan}
The following workflow plan can be used as a contributor's guide to the development process. 

\textbf{Issues and deliverables management}:
Each deliverable (and sub-deliverable if needed) will have a milestone associated with it. 
Each milestone will have a set of issues associated with it, representing the development tasks needed to complete the deliverable.
These issues will be in the project Kanban board, and will be moved through the columns (To Do, In Progress, Done) as they are worked on. \newline

\noindent \textbf{Pull requests and code review}:
All code changes will be made in a new branch, and a pull request will be created once development is done for a particular task. See the branches and commit sections for naming convention. 
Typically, a branch will be created for each issue, but a branch may close several issues if they are related and can be completed together.
Pull requests will be reviewed by at least one other team member, and the pull request reviewer role will be responsible for approving or requesting changes to the pull request.

\subsection{Branches \& Commits}
\textbf{Branch naming convention:} Each branch will be named according to the following convention: \texttt{<branch-type>/<branch-name>}.
Like mentioned, development will be done in a new branch, and a pull request will be created once development is done for a particular task.
The naming convention for branches will be kebab-case and will be based on the branch content as well as the branch type. 
For example, a branch for the issue "Implement user login" would be named "feature/implement-user-login".
The current types of branches are feature, docs, and bugfix, but more types may be added as needed. \newline
\textbf{Commit naming convention:} Each commit will be named according to the following convention: \texttt{<commit-type>: <commit-message>}.
The naming convention for commits will be lowercase and will be based on the commit content as well as the commit type.
For example, a commit for the issue "Implement user login" would be named "feat: implement user login".
The current types of commits are feat, fix, docs, style, refactor, and test, but more types may be added as needed. \newline

\subsection{Merging strategy}
Developers will merge their changes into the main branch through pull requests.
A pull request must be approved by at least one other team member before it can be merged into the main branch.
Individual commits will be preserved in the pull request, and the pull request reviewer will be responsible for approving or requesting changes to the pull request.

\section{Project Decomposition and Scheduling}

The project is organized around key milestones and overarching stages 
(refer to Table 2), progressing from initial project planning and requirements
analysis through design, implementation, verification, and final demonstration. 
Each stage contains defined deliverables and deadlines to provide the team 
with a clear progression throughout the project. The schedule will remain 
flexible where necessary, allowing the team to adjust tasks and responsibilities 
based on project needs, progress, and individual strengths.

\begin{table}[ht]
    \centering
    
    \begin{tabularx}{\textwidth}{|>{\centering\arraybackslash}p{3cm}|
                                    X|
                                    >{\centering\arraybackslash}p{2cm}|}
        \hline
        \textbf{Stage} & \textbf{Milestone} & \textbf{Deadline} \\
        \hline
        
        Planning & Project selection & W2 \\
        \hline
        
        Planning & Problem Statement + PoC + Development Plan + Team Charter & W4 \\
        \hline
        
        Requirements & SRS + Hazard Analysis & W6 \\
        \hline
        
        Verification & V\&V Plan & W8 \\
        \hline
        
        Design & Design Document & W10 \\
        \hline
        
        Design & PoC Demo & W11 \\
        \hline
        
        Implementation & Design Document Rev 0 & W16 \\
        \hline
        
        Implementation & Module Implementation & W16 \\
        \hline
        
        Implementation & Rev 0 Demo & W18 \\
        \hline
        
        Finalization & Final Demo + Documentation & W21 \\
        \hline
        
        Finalization & V\&V Report + CD & W22 \\
        \hline
        
        Finalization & Final Demos & W23 \\
        \hline
        
        Expo & Expo Materials & W25 \\
        \hline
        
        Expo & Capstone EXPO & W26 \\
        \hline
    \end{tabularx}
    
    \caption{Project Milestones and Schedule}
    \label{tab:milestones}
\end{table}

\subsection{GitHub Project Usage}
The project board will follow a Kanban-style workflow with columns for Backlog,
Ready, In Progress, In Review, and Done. The project will be organized into 
milestones representing the overarching deliverables required throughout each 
stage of the project. For example, the initial stage contains milestones for 
the Project Statement and Development Plan, with individual tasks associated 
with each milestone. Each task will be represented as a card on the board and 
assigned to the appropriate team member. Cards will move through the Kanban columns 
as work progresses, allowing the team to track the status of individual tasks and overall 
milestone completion.


\section{Proof of Concept Demonstration Plan}

\subsection{Project Risks}

One main risk is designing an aggregation algorithm that combines predictions
from multiple inference providers into an accurate final result. The algorithm
must handle disagreements between providers and reduce the influence of
dishonest or low-quality predictions. This is difficult because, for a given
task, the network may not have a source of \textit{ground truth} against which
to compare the predictions. There is also a cold-start problem where intially
there may not be enough trusted providers in the network to correctly evaluate
one another. 

Another risk is designing the incentive system and token economics so that the network
is sufficiently decentralizated and is resistant to attacks. The system must be fair and
and accesible for participants so that they are encouraged to take part in the network.
The system must also be economically sustainable and ensure that
the distribution of rewards does not create incentives to exploit the system
or reduce the quality of the aggregated results.

Collusion attacks pose a significant risk to the system. A collusion attack
occurs when multiple participants, such as inference providers, coordinate
their actions to manipulate the aggregation algorithm, reputations, and reward
distribution. For example, providers could submit similar inaccurate
predictions to make their results appear correct and receive rewards. This is
difficult to detect when there is no ground truth and when the colluding
providers make up a large portion of the network.

\subsection{POC Demonstration Plan}

The proof of concept plan is to create a simulation of the DAIIM network that
is completely off chain. The simulation will execute a sequence of rounds in which
an inference task for a regression problem is executed by the network. Various different types of actors will 
participate in the network:

1. Honest: attempt to maximize rewards by providing accurate predictions consistently (with some noise)
2. Learners: gradually improve their predictions over time
3. Malicious: intentiononally provide bad predictions
4. Collusion: A group of providers working together to skew the network

The simulation will start with a simple ML task (eg. predict y, y = mx + b), and each round some input data will be randomly generated.
The different actors will be given some algorithm (good or bad algorithm depending on the actor) to determine their predictions given the input.
The actors will submit their predictions, and the aggregation algorithm will determine the aggregated result. Actor's reputations will be updated
and rewards will be distributed after each round.

Another layer of this simulation is to economically model each actor. We will
assign providers a cost for running their model, transaction fees for submitting
predictions on-chain, and economic rewards or penalties for participating. These
values will be calibrated relative to the costs and payments that we expect a
provider to realistically encounter when performing an inference task. This
will ensure that the simulation reflects realistic participation decisions
rather than relying on arbitrary values. We will also vary these values to
measure how sensitive participation and network performance are to changes in
costs, fees, rewards, and penalties.

We will test various combinations of actors to evaluate the limits of the
system. These scenarios will vary the proportions of different actors, as well as the number of providers in the
network. These results will allow us to identify conditions under which the system remains reliable and economically
viable, as well as conditions under which malicious or colluding providers can
significantly influence the network.

\subsection{Success Criteria and Failure Response}

This proof of concept will be successful and show that the stated risk can be overcome if:
- The system provides accurate aggregated results that improve over time even with some limited number of malicious actors
- It is more profitable to act honestly and attempt to provide good predictions than to be malicious
- The network can be cold-started with the assumption that initially there are a majority of honest providers

If the proof of concept is not able to demonstrate that these risks can be addressed then the project's assumptions and scope
may need to be changed. Before proceeding to full
implementation, the team will review the findings with the supervisor and
decide whether to revise the design, broaden the threat
conditions, or reconsider the project's feasibility.

\section{Expected Technology}

This section describes the technologies that are expected to support the
implementation of DAIIM. Note that these choices are preliminary and subject to change later in the project.

\subsection{Core Implementation Stack}

DAIIM will use Solana as the blockchain network for recording transactions and
executing on-chain logic. During development, the team will use Solana Devnet,
a separate network for testing programs and transactions without affecting
Solana Mainnet. We will request free test SOL from a Devnet faucet to pay for
transaction fees. Since Devnet tokens have no real-world value, the team can
develop and test DAIIM without using real money before deploying to Mainnet.

Solana programs, or smart contracts, will be written in Rust. The conventional
off-chain backend will use Node.js with TypeScript and PostgreSQL for data
storage. The user-facing interface will be developed with React.js and Vite.
Python may also be used to write simple machine learning models for integration
testing purposes.

\subsection{Libraries and External Components}

The expected external libraries and components are:

\begin{itemize}
  \item \textit{[Anchor]}: A development framework for building, deploying and testing Solana programs in Rust.
  \item \textit{[Metamask]}: Used for accessing crypto wallets \& signing transactions.
  \item \textit{[Metaplex]}: Protocol and SDK for creating NFTs.
  \item \textit{[@solana/web3.js]}: Javascript library for interacting with the Solana blockchain.
  \item \textit{[@solana/spl-token]}: JavaScript library for interaction with the SPL token program on Solana.
  \item \textit{[Vitest]}: Testing framework for the React frontend and other TypeScript code.
  \item \textit{[Solana Devnet JSON-RPC]}: Public RPC access for testing blockchain transactions and programs.
\end{itemize}

\subsection{Version Control and CI/CD}
Git and Github will be used for version control and the project repository. GitHub Projects will
organize issues and track their progress. LaTeX will be used for the
project's technical documentation.

GitHub Actions will compile the LaTeX documentation automatically when relevant
changes are pushed or proposed. Later, the CI pipeline will be extended to
build and run tests for the solana program (using Anchor) and React frontend.

\section{Use of AI and LLMs}

\wss{This section is a continuation of the previous section.  The role of LLMs
in your project as important enough to warrant their own section.  
How will your team be using LLMs for code and documentation development?  
What tools will you be using?  How will you verify the output of your LLMs?}

\section{Coding Standard}

The team will adhere to the coding standards outlined in Table~\ref{tab:coding-standards}.

\begin{table}[ht]
    \centering
    
    \begin{tabularx}{\textwidth}{|>{\centering\arraybackslash}p{3cm}|X|}
        \hline
        \textbf{Category} & \textbf{Details} \\
        \hline
        
        \textbf{React Frontend} &
        React code will follow the
        \href{https://react.dev/learn}{official React documentation}
        and established JavaScript/TypeScript conventions. Code will be
        formatted using \texttt{Prettier} and linted using \texttt{ESLint}. \\
        \hline

        
        \textbf{Rust Backend} &
        Rust code will follow the official
        \href{https://rust-lang.github.io/api-guidelines/}{Rust API Guidelines}
        and community style conventions, formatted using \texttt{rustfmt}
        and linted with \texttt{clippy}. \\
        \hline
        
        \textbf{Rust Documentation} &
        Rust functions, structs, and modules will use
        \texttt{rustdoc} doc comments (\texttt{///} and \texttt{//!})
        following the conventions in the official
        \href{https://doc.rust-lang.org/rustdoc/how-to-write-documentation.html}
        {Rust documentation guide}. \\
        \hline
        
        \textbf{Docker} &
        Dockerfiles and container images will follow
        \href{https://github.com/microsoft/docker/blob/master/docs/userguide/eng-image/dockerfile_best-practices.md}
        {Docker's official best practices}
        for maintainability, efficiency, and security. \\
        \hline
        
        \textbf{Git} &
        Git will be used for version control, with feature branches
        used to isolate development work. \\
        \hline
        
        \textbf{Code Review} &
        Changes will be submitted through GitHub pull requests and
        reviewed by another team member before being merged into the
        main branch. See the
        \href{https://docs.github.com/en/pull-requests/concepts/giving-reviews}
        {GitHub pull request review guidelines}. \\
        \hline
        
    \end{tabularx}
    
    \caption{Coding Standards}
    \label{tab:coding-standards}
    
\end{table}

\newpage{}

\section{Appendix --- Generative AI Use Disclosure}

\wss{ChatGPT (version 5) was used to brainstorm ideas for the template for this disclose section.}

\wss{This section is about the use of AI to write this document, not your
planned use of AI for your project.  You discuss that topic in one of the
sections in the body of this report.}

\subsection{AI Tools Used}

\wss{Give the name of the AI tool(s) used and the version.}

\subsection{How and Where Used}

\wss{Explain what the tool(s) were used for.  Examples include: brainstorming,
explaining a technical concept, reviewing the document, generating or modifying
text, generating or modifying diagrams, identifying errors or omissions or
inconsistencies.}

\wss{For each of the uses, identify which parts of the document were affected.}

\wss{You do not need to report every interaction, but you should disclose the 
uses that materially contributed to your work}

\subsection{Verification of AI-Generated Content}

\wss{\begin{itemize}
    \item Describe how the team checked AI-generated or AI-assisted material
    before including it in the document.
    \item In particular, verify factual claims, technical assertions,
    requirements, calculations, references, and other information for which an
    AI system may produce incorrect or fabricated results.
    \item \textbf{The team is responsible for the accuracy, completeness,
consistency, and appropriateness of the submitted document, regardless of
whether material was generated by a human or an AI system.}
\end{itemize}}

\newpage{}

\section*{Appendix --- Reflection}

\wss{Not required for CAS 741}

\input{../Reflection.text}

\begin{enumerate}
    \item Why is it important to create a development plan prior to starting the
    project?
    \item In your opinion, what are the advantages and disadvantages of using
    CI/CD?
    \item What disagreements did your group have in this deliverable, if any,
    and how did you resolve them?
\end{enumerate}

\newpage{}

\section*{Appendix --- Team Charter}
\section*{External Goals}

For this Capstone, our team hopes to deliver a project that is both challenging and enjoyable to create. We aim to build something beyond simply fulfilling the requirements of a school project---something that we can be proud of and potentially carry forward into our professional careers. By focusing on producing a meaningful, well-developed, and high-quality project, we believe that a strong performance in the Capstone course will naturally follow.

\section*{Attendance}

This section summarizes the ground rules and expectations regarding team member attendance.

\subsection*{Expectations}

All team members are expected to attend weekly team meetings and meetings with our supervisor. Attendance at every Capstone course lecture is not required, provided that members remain up to date with the course material through either attending the lecture in person or reviewing the material independently.

\vspace{\baselineskip}
\noindent
Attendance at all required team and supervisor meetings will be tracked through GitHub Issues. All members are expected to maintain a 95\% attendance rate throughout the project.
\vspace{\baselineskip}
\noindent
When a meeting is held in person, team members should make their best effort to attend in person. If a member is unable to attend in person, they should notify the team at least 5 hours in advance and attend virtually instead. If a member is unable to attend the meeting altogether, they must notify the team as soon as possible, preferably when they become aware of the conflict and at least 24 hours before the meeting when reasonably possible. Different expectations may apply depending on the circumstances.

\subsection*{Acceptable Excuse}

Team members are expected to communicate with the rest of the team if they are unable to attend a meeting or meet a project deadline. Each member may miss up to three team meetings per term with prior notice. Acceptable reasons may include illness, family emergencies, religious observances, or previously arranged commitments that cannot reasonably be rescheduled.
\vspace{\baselineskip}

\noindent
If a member exceeds three absences in a term, the team will discuss the member's participation, expectations, and responsibilities. If the issue cannot be resolved within the team, the team may involve the course professor or supervisor.
\vspace{\baselineskip}

\noindent
Team members are expected to meet all assigned deadlines. Exceptions will generally only be made in the event of an emergency, as outlined in the following section.

\subsection*{In Case of Emergency}

If a team member is unable to meet a deadline or complete their assigned work due to an emergency, they must communicate this to the team as soon as they become aware that the work cannot be completed on time.
\vspace{\baselineskip}
\noindent
The team will then discuss appropriate accommodations and, if necessary, redistribute the affected work to ensure that project deadlines can still be met. If another team member must take on the incomplete work, the original team member will be expected to take on additional responsibilities during a subsequent phase, where appropriate, to help maintain a fair distribution of workload.


\wss{borrows from
\href{https://engineering.up.edu/industry_partnerships/files/team-charter.pdf}
{University of Portland Team Charter}}


\subsection*{Accountability and Teamwork}

\subsubsection*{Quality} 

\wss{What are your team's expectations regarding the quality
of team members' preparation for team meetings and the quality of the
deliverables that members bring to the team?}

\subsubsection*{Attitude}

\wss{What are your team's expectations regarding team members' ideas,
interactions with the team, cooperation, attitudes, and anything else regarding
team member contributions?  Do you want to introduce a code of conduct?  Do you
want a conflict resolution plan?  Can adopt existing codes of conduct.}

\subsubsection*{Stay on Track}

\wss{What methods will be used to keep the team on track? How will your team
ensure that members contribute as expected to the team and that the team
performs as expected? How will your team reward members who do well and manage
members whose performance is below expectations?  What are the consequences for
someone not contributing their fair share?}

\wss{You may wish to use the project management metrics collected for the TA and
instructor for this.}

\wss{You can set target metrics for attendance, commits, etc.  What are the
consequences if someone doesn't hit their targets?  Do they need to bring the
coffee to the next team meeting?  Does the team need to make an appointment with
their TA, or the instructor?  Are there incentives for reaching targets early?}

\subsubsection*{Team Building}

\wss{How will you build team cohesion (fun time, group rituals, etc.)? }

\subsubsection*{Decision Making} 

\wss{How will you make decisions in your group? Consensus?  Vote? How will you
handle disagreements? }

\end{document}